% Chunk: Chapter 11 opening and \S11.1, Counterterms.
% Source: Weinberg QFT Vol. I, physical pp. 495--496 (printed pp. 472--473),
% ending immediately before the \S11.2 heading on physical p. 496.
% Coverage: equations (11.1.1)--(11.1.9); one footnote; no figures.
% Uncertainty: none.

In this chapter we shall proceed to carry out some of the classic one-loop
calculations in the theory of charged leptons---massive spin $\tfrac12$ particles
that interact only with the electromagnetic field. There are three known
species or ``flavors'' of leptons: the electron and muon, and the heavier,
more recently discovered tauon. For definiteness we shall refer to the
charged particles in our calculations here as ``electrons,'' though most of
our calculations will apply equally to muons and tauons. After some
generalities in Section 11.1, we will move on to the calculation of the
vacuum polarization in Section 11.2, the anomalous magnetic moment of
the electron in Section 11.3, and the electron self-energy in Section 11.4.
Along the way, we will introduce a number of the mathematical techniques
that prove useful in such calculations, including the use of Feynman
parameters, Wick rotation, and both the dimensional regularization of
't Hooft and Veltman and the older regularization method of Pauli and
Villars. Although we shall encounter infinities, it will be seen that the
final results are finite if expressed in terms of the renormalized charge
and mass. In the next chapter we shall extend what we have learned
here about renormalization to general theories in arbitrary orders of
perturbation theory.

\subsection{Counterterms}

The Lagrangian density for electrons and photons is taken in the form\footnote{In
this chapter we will not be making transformations between Heisenberg- and
interaction-picture operators, so we shall return to a conventional notation,
in which an upper case $A$ and a lower case $\psi$ are used to denote the
photon and charged particle fields, respectively.}
\begin{equation}
  \mathcal L=-\frac14 F_{\mathrm B}^{\mu\nu}F_{\mathrm B\,\mu\nu}
  -\bar\psi_{\mathrm B}\bigl[\gamma_\mu(\partial^\mu+\ii e_{\mathrm B}A_{\mathrm B}^{\mu})
  +m_{\mathrm B}\bigr]\psi_{\mathrm B},
  \tag{11.1.1}\label{eq:11.1.1}
\end{equation}
where $F_{\mathrm B}^{\mu\nu}\equiv\partial^\mu A_{\mathrm B}^{\nu}
-\partial^\nu A_{\mathrm B}^{\mu}$; $A_{\mathrm B}^{\mu}$ and
$\psi_{\mathrm B}$ are the bare (i.e., unrenormalized) fields of the photon
and electron, and $-e_{\mathrm B}$ and $m_{\mathrm B}$ are the bare charge and
mass of the electron. As described in the previous chapter, we introduce
renormalized fields and charge and mass:
\begin{align}
  \psi&\equiv Z_2^{-1/2}\psi_{\mathrm B},
  \tag{11.1.2}\label{eq:11.1.2}\\
  A^\mu&\equiv Z_3^{-1/2}A_{\mathrm B}^{\mu},
  \tag{11.1.3}\label{eq:11.1.3}\\
  e&\equiv Z_3^{+1/2}e_{\mathrm B},
  \tag{11.1.4}\label{eq:11.1.4}\\
  m&\equiv m_{\mathrm B}+\delta m,
  \tag{11.1.5}\label{eq:11.1.5}
\end{align}
with the constants $Z_2$, $Z_3$, and $\delta m$ adjusted so that the
propagators of the renormalized fields have poles in the same position and
with the same residues as the propagators of the free fields in the absence
of interactions. The Lagrangian may then be written in terms of renormalized
quantities, as
\begin{equation}
  \mathcal L=\mathcal L_0+\mathcal L_1+\mathcal L_2,
  \tag{11.1.6}\label{eq:11.1.6}
\end{equation}
where
\begin{equation}
  \mathcal L_0=-\frac14 F^{\mu\nu}F_{\mu\nu}
  -\bar\psi\bigl[\gamma_\mu\partial^\mu+m\bigr]\psi,
  \tag{11.1.7}\label{eq:11.1.7}
\end{equation}
\begin{equation}
  \mathcal L_1=-\ii e A_\mu\bar\psi\gamma^\mu\psi,
  \tag{11.1.8}\label{eq:11.1.8}
\end{equation}
and $\mathcal L_2$ is a sum of ``counterterms''
\begin{align}
  \mathcal L_2={}&-\frac14(Z_3-1)F^{\mu\nu}F_{\mu\nu}
  -(Z_2-1)\bar\psi\bigl[\gamma_\mu\partial^\mu+m\bigr]\psi
  \notag\\
  &+Z_2\delta m\,\bar\psi\psi
  -\ii e(Z_2-1)A_\mu\bar\psi\gamma^\mu\psi.
  \tag{11.1.9}\label{eq:11.1.9}
\end{align}
It will turn out that all of the terms in $\mathcal L_2$ are of second order
and higher order in $e$, and that these terms just suffice to cancel the
ultraviolet divergences that arise from loop graphs.
